\setcounter{section}{1}
\setcounter{theorem}{7}
\section{Lower bounds for constant depth linear circuits}\label{sec:shoup-smol based lb}
In this section, we prove Theorem 2 of the original publication. 
We start by describing  the notion of Shoup-Smolensky dimension, but first we set up some notation.

\subsection{Notation}
 We work with matrices whose entries lie in an appropriate extension of a base finite field $\F_p$. We follow the natural convention that the elements of this extension will be represented as univariate polynomials of appropriate degree over the base field, and the arithmetic is done modulo an explicitly given irreducible polynomial. 

We use boldface letters ($\vecx,\vecy$) to denote vectors. The length of the vectors is understood from the context.

For a matrix $M$, $\spars{M}$ denotes the number of non-zero entries in $M$. 

\subsection{Shoup-Smolensky Dimension}
\label{sec:ssdim}
A useful concept will be the notion of Shoup-Smolensky dimension of sequences of elements of an extension $\E$ of a field $\F$. 
\begin{definition}[Shoup-Smolensky dimension]\label{def:SS-dim}
Let $\F$ be a field, and $\E$ be an extension field of $\F$. Let $S=(a_1,\ldots,a_m)$ a sequence of elements of $\E$. For $t\in\N$, denote by $\Pi_t(S)$ the set of $t$-wise products of distinct entries of $S$ that is,
\[\Pi_t(S)=\left\{\prod_{j=1}^t a_{i_j}:1\le i_1<i_2<\cdots<i_t\le m\right\}.\]
The Shoup-Smolensky dimension of $S$ of order $t$, denoted by $\Gamma_{t,\F}(S)$ is defined to be the dimension, over $\F$, of the vector space spanned by $\Pi_t(S)$.
We also denote by $\Sigma_t(S)$ the number of distinct elements of $\E$ that can be obtained by summing distinct elements of $\Pi_t(S)$.
When $M\in\E^{n\times n}$ is a matrix we also regard it as a sequence of $m=n^2$ elements of $\E$ (under some order on the entries) and refer to the Shoup-Smolensky dimension of $M$.
\end{definition}


\subsection{Upper bounding the Shoup-Smolensky dimension for Sparse Products}

\label{sec:SS-dim}
The following lemma shows that any matrix computable by a depth-$d$ linear circuit of size at most $s$ has a somewhat small Shoup-Smolensky dimension. 

\begin{lemma}\label{lem:ss-easy-ub-large-depth}
Let $\F$ be a field, $\E$ an extension of $\F$ and $A \in \E^{n \times n}$ be a matrix such that $A = \prod_{i = 1}^d P_i$ for $P_i \in \E^{n_i \times m_i}$, where $\sum_{i = 1}^d\spars{P_i} \leq s$. Then, for every $t\le n^2/4$ such that $s \ge dt$ it holds that
\[
\Gamma_{t, \F}(A) \le \inparen{e^d (2s/dt)^d}^t.
\]
\end{lemma}
\begin{proof}
Since 
$$A_{i,j} = \left(\prod_{\ell = 1}^d P_{\ell}\right)_{i,j} = \sum_{k_1, \ldots, k_{d-1}} (P_1)_{i, k_1}\cdot \left(\prod_{\ell = 2}^{d-1} (P_{\ell})_{k_{\ell-1}, k_{\ell}} \right) \cdot (P_{d})_{k_{d-1}, j}\, , $$
every element in $\Pi_{t} (A)$ is a sum of monomials of degree $dt$ in the entries of $P_1, P_2, \ldots, P_d$, that is,
\[
\Gamma_{t, \F}\left(\prod_{i = 1}^d P_i\right) \leq \binom{s + dt}{dt},
\]
with the right hand side being the number of monomials of degree $dt$ in $s$ variables.
Using the inequality $\binom{n}{k} \le (en/k)^k$,
\[
\Gamma_{t, \F}(A) \leq (e(1 + s/dt))^{dt} \le \inparen{e^d (2s/dt)^d}^t. \qedhere
\]
\end{proof}


Over $\Q$, we do not wish to use field extensions (which would give rise to elements with infinite bit complexity). Thus, we use a similar argument that replaces the measure $\Gamma_{t,\F}$ with $\Sigma_{t}$ (recall \autoref{def:SS-dim}) for a small tolerable penalty.
\begin{lemma}
\label{lem:ss-up-sigma}
Let $d$ be a positive integer. Let $A \in \Q^{n \times n}$ be a matrix such that $A = \prod_{i = 1}^d P_i$ for $P_i \in \Q^{n_i \times m_i}$, where $\sum_{i = 1}^d\spars{P_i} \leq s $. Assume that for each $i$,  $n_i \leq n^2$ and $m_i \leq n^2$. Then, for every $t\le n^2/4$ such that $s \ge dt$ it holds that
\[
\Sigma_{t}(A) \le 2^{2n^3\cdot \inparen{e^d (2s/dt)^d}^t}.
\]

\end{lemma}

\begin{proof}
We follow the same steps as in the proof of~\autoref{lem:ss-easy-ub-large-depth}, replacing the measure $\Gamma_{t,\F}(A)$ by $\Sigma_t(A)$. As before, 
\[
A_{i,j} = \left(\prod_{\ell = 1}^d P_{\ell}\right)_{i,j} = \sum_{k_1, \ldots, k_{d-1}} (P_1)_{i, k_1}\cdot \left(\prod_{\ell = 2}^{d-1} (P_{\ell})_{k_{\ell-1}, k_{\ell}} \right) \cdot (P_{d})_{k_{d-1}, j}\, .
\]
Every element in $\Pi_{t} (A)$ can be written as 
\begin{equation}
\label{eq:monomials-in-product}
\sum_{\alpha \in \mathcal{M}} c_\alpha \cdot \alpha
\end{equation}
where $\mathcal{M}$ is the set of monomials of degree $dt$ in the entries of $P_1, P_2, \ldots, P_d$, and each $c_\alpha$ is a non-negative integer of of absolute value at most $s^{dt} \le 2^{n^3}$ (since $s \leq n^2d$ and $d$ is $O(1)$). It now follows that each element in $\Sigma_t(A)$ has the same form as in \eqref{eq:monomials-in-product}, with $c_\alpha \le |\Pi_t(A)| \cdot 2^{n^3} \le 2^{2n^3}$ . We conclude that
\[
\Sigma_t(A) \le (2^{2n^3})^{\binom{s+dt}{dt}},
\]
which implies the statement of the lemma using the same bounds on binomial coefficients as in \autoref{lem:ss-easy-ub-large-depth}.
\end{proof}


We now move on to describe constructions of matrices which have large Shoup-Smolensky dimension, and then deduce lower bounds for them.

\subsection{Sidon sets and hard univariate matrices}\label{sec:sidon-ss dim}
In this section, we describe a construction of a matrix $G \in \F[y]^{n \times n}$ which has a large value of $\Gamma_{t, \F}$. Let us denote $G_{i,j} = y^{e_{i,j}}$ for some non-negative integer $e_{i,j}$. For $G$ to have a large Shoup-Smolensky dimension of order $t$, the set $S = \set{e_{1,1}, e_{1,2}, \ldots, e_{n,n}} \subseteq \N$ should have the property that $S^{(t)} := \set{a_1 + a_2 + \ldots + a_t : a_i \in S \text{ distinct}}$ has size comparable to $\binom{|S|}{t}$. A set $S$ such that every subset of size $t$ of $S$ has a distinct sum is called a \emph{$t$-wise Sidon set}. These are very well studied objects in arithmetic combinatorics, and explicit constructions are known for them in $\poly(n)$ time (e.g., Lemma 60 in~\cite{Bshouty}). However, another important parameter in the construction is the degree of $y$, and such a set will inevitably contain integers of size roughly $n^{\Omega(t)}$. Thus, the construction of $G$ would take time which is not polynomially bounded in $n$. Below we give an elementary construction of such a set in time $n^{O(t)}$ (cf.\ \cite{AGKS15}).

\begin{lemma}\label{lem:easy-sidon}
Let $t$ be a positive integer. There is a set $S\subseteq\N$ of size $m$ such that:
\begin{enumerate}
\item $S^{(t)}:=\{a_1+a_2+\cdots+a_t:a_i\in S\text{ distinct}\}$ has size $\binom mt$.
\item The maximal element in $S$ is at most $m^{O(t)}$.
\item $S$ can be constructed in time $m^{O(t)}$.
\end{enumerate}
\end{lemma}
\begin{proof}
Let $S'=\{1,2,2^2,\ldots,2^{m-1}\}$. Clearly, every subset of $S'$ has a distinct sum. For a prime $p$ we denote $S_p=S'\bmod p=\{a\bmod p:a\in S'\}$, and we claim that there exists a prime $p\le m^{O(t)}$ such that $|(S_p)^{(t)}|=\binom mt$. Since this condition can be checked in time $m^{O(t)}$, this would immediately imply the statement of the lemma, by checking this condition for every $p\le m^{O(t)}$ and letting $S=S_p$ for a $p$ which satisfies this condition.
For every subset $T\subseteq S'$ of size $t$, let $\sigma_T$ denote the sum of its elements, and observe that $\sigma_T\le2^m$. Clearly, $\sigma_T\bmod p=\sigma_{T'}\bmod p$ if and only if $p\mid\sigma_T-\sigma_{T'}$, so it is enough to show that there exists $p\le m^{O(t)}$ which does not divide
\[N:=\prod_{\substack{T\ne T'\subseteq S'\\|T|=|T'|=t}}(\sigma_T-\sigma_{T'}),\]
and therefore does not divide any of the terms on the right hand size. It further holds that $0\ne N\le(2^m)^{m^{O(t)}}=2^{m^{O(t)}}$, so the existence of $p$ now follows from the fact that $N$ can have at most $\log N=m^{O(t)}$ distinct prime divisors, and from the prime number theorem.
\end{proof}

Given the above construction of $t$-wise Sidon sets, we now describe the construction of matrices with univariate polynomial entries which has large Shoup-Smolensky dimension.

\begin{construction}
\label{con:univariate hard matrix}
Let $S = \set{e_{i, j} : i, j \in [n]}$ be a $t$-wise Sidon set of positive integers of size $n^2$, as in \autoref{lem:easy-sidon}. Then, the matrix $G_{t,n} \in \F[y]^{n \times n}$ is defined as follows as 
$(G_t)_{i, j} = y^{e_{i, j}}$.
\end{construction}

The useful properties of \autoref{con:univariate hard matrix} are given by the following lemma.

\begin{lemma}
\label{lem:univariate hard properties}
Let $t \le n$ be a parameter, $S\subseteq\N$ be a $t$-wise Sidon set of size $n^2$ and let $G_{t,n}$ be the matrix defined in~\autoref{con:univariate hard matrix}. Then, the following are true. 
\begin{enumerate}
\item Every entry of $G_{t,n}$ is a monomial of degree at most $n^{O(t)}$. 
\item $\Gamma_{t,\F}((G_{t,n})) \geq \binom{n^2}{t} \ge  \left(\frac{n^2}{t}\right)^t$.
\end{enumerate}
\end{lemma}
\begin{proof}
The first item follows from the definition of $G_{t,n}$ and the properties of the set $S$ in~\autoref{lem:easy-sidon}. The second item also follows from the properties of $S$ and the definition of Shoup-Smolensky dimension, since every $t$-wise product of elements of $G_{t,n}$ gives a distinct monomial in $y$, and thus they are all linearly independent over the base field $\F$.
\end{proof}

\subsection{Hard matrices over finite fields}
From the univariate matrix in \autoref{con:univariate hard matrix}, we  now construct, for every $p$ and parameter $t$, a matrix $M$ over an extension of $\F_p$ which has large Shoup-Smolensky dimension over $\F_p$ with the same parameters as $G_{t,n}$.
 
\begin{lemma}
\label{lem:hard-over-finite}
Let $p$ be a prime, and $t $ be any positive integer. There is a matrix $M_{t,n} \in \E^{n\times n}$ over an extension $\E$ of $\F_p$ of degree $\exp\inparen{{O(t\log n)}}$, which can be deterministically constructed in time $n^{O(t)}$, and satisfies
\[
\Gamma_{t, \F_p}(M_{t,n}) \geq \left(\frac{n^2}{t}\right)^t\, 
\] 
\end{lemma}
\begin{proof}
Let $G_{t,n}$ be as in \autoref{con:univariate hard matrix}, and let $\Delta$ be the maximum degree of any entry of $G_{t,n}$. Set $D = 10\cdot t\cdot \Delta = \exp\left(O(t\log n)\right)$. We use Shoup's algorithm (see Theorem 3.2 in~\cite{Shoup90}) to construct an irreducible polynomial $g(z)$ of degree $D+1$ over $\F_p$ in deterministic  $\poly(D, |\F_p|)$ time. Let $\alpha$ be a root of $g(z)$ in an extension $\E$ of $\F_p$, where $\E \equiv \F_p[z]/\langle g(z) \rangle$.\footnote{We identify the elements of $\E$ with coefficient vectors of polynomials of degree at most $D$ in $\F_p[z]$, and in this representation $\alpha$ is identified with the polynomial $z$.} Then, it follows that $1, \alpha, \alpha^2, \ldots, \alpha^{D}$ are linearly independent over $\F$.

The matrix $M_{t,n}$ is obtained from $G_t$ by just replacing every occurrence of the variable $y$ by $\alpha$. We now need to argue that $M_{t,n}$ continues to satisfy $\Gamma_{t, \F_p}(M_{t,n}) \geq \left(\frac{n^2}{t}\right)^t$. By the choice of $\alpha$, it immediately follows that $\Gamma_{t, \F_p}(M_{t,n}) = \Gamma_{t, \F_p}(G_{t,n})$, since every monomial in the set $\Pi_t(M_{t,n})$ is mapped to a distinct power of $\alpha$ in $\{0, 1, \ldots, D\}$, which are all linearly independent over $\F_p$.

The upper bound on the running time needed to construct $M_{t,n}$ now follows from the upper bound on the degree of the extension $\E$, and from \autoref{lem:easy-sidon}.
\end{proof}
The following theorem now directly follows.
\begin{theorem}\label{thm:depth-d-finite}
Let $p$ be any prime and $d\geq 2$ be a positive integer. Then, there exists a family of matrices $\{A_n\}_{n \in \N}$ which can be constructed in time $n^{O(n^{1-1/2d})}$ such that every depth-$d$ linear circuit $\overline{\F}_p$ computing $A_n$ has size at least $\Omega(n^{1 + 1/2d})$. Moreover, the entries of $A_n$ lie in an extension of $\F_p$ of degree at most $\exp(O(n^{1-1/2d}\log n))$. 
\end{theorem}
\begin{proof}
We invoke~\autoref{lem:hard-over-finite} with parameter $t$ set to  $n^{1-1/2d}$ to get matrices $\{A_n\}$ in time $n^{O(t)}$ with the following lower bound on their Shoup-Smolensky dimension. 
\[
\Gamma_{t, \F_p}(M_{n}) \geq \left(\frac{n^2}{t}\right)^t\, .
\] 
If there is a  depth $d$ linear circuit of size $s$ computing the linear transformation $A_n\cdot \vecx$, the following inequality must hold (from~\autoref{lem:ss-easy-ub-large-depth}),
\begin{equation}\label{eqn:depth-d-lb}
\inparen{e^d (2s/dt)^d}^t \geq \left(\frac{n^2}{t}\right)^t \, .
\end{equation}
If $s \leq   n^{1+1/2d}/2$, we have, 
\[
\inparen{e^d (2s/dt)^d}^t \leq (O(e/d))^{dt} \cdot n^t\, . 
\]
We also have, 
\[
\left(\frac{n^2}{t}\right)^t \geq \left(n^{1 + 1/2d} \right)^t \, .
\]
For any constant $d$, these estimates contradict~\autoref{eqn:depth-d-lb}, thereby implying a lower bound of $\Omega(n^{1+1/2d})$ on s. 
\end{proof}

